\documentclass[11pt]{article}

\usepackage[T1]{fontenc}
\usepackage[utf8]{inputenc}
\usepackage{lmodern}
\usepackage[margin=1in]{geometry}
\usepackage{amsmath,amssymb,amsthm,mathtools}
\usepackage{hyperref}
\usepackage{microtype}

\hypersetup{colorlinks=true,linkcolor=blue,urlcolor=blue,citecolor=blue}

\newtheorem{theorem}{Theorem}
\newtheorem{corollary}[theorem]{Corollary}
\newtheorem{proposition}[theorem]{Proposition}
\newtheorem{remark}[theorem]{Remark}

\newcommand{\dd}{\,\mathrm{d}}
\newcommand{\RR}{\mathbb{R}}
\newcommand{\ZZ}{\mathbb{Z}}
\newcommand{\GG}{\Gamma}
\newcommand{\E}{\mathcal{E}}
\newcommand{\mean}[1]{\overline{#1}}

\title{An Exact Minimum for Weighted Closed-Loop Phase Locking}
\author{}
\date{}

\begin{document}
\maketitle

\begin{abstract}
For a prescribed closed loop, positive weight, geometric phase target, connection,
and winding number, we determine the exact minimum of a weighted quadratic
phase-mismatch functional.  A sharp Cauchy--Schwarz inequality, together with
an explicit admissible phase, gives the minimum and a unique minimizing
covariant phase gradient.  An exact remainder identity establishes strict
minimality modulo constant-phase symmetry.  Comparison between different
winding sectors is conditional on the locking functional accounting for all
sector-dependent contributions and does not itself provide a transition
mechanism.
\end{abstract}

\section{Setting and admissible phases}

Let $\GG$ be an oriented closed loop of total arclength $S>0$, with coordinate
$s\in[0,S]$ and its endpoints identified.  Fix integers $m,n\in\ZZ$.
Let $w:[0,S]\to\RR$ be measurable, and suppose there are constants
$0<w_-\leq w_+<\infty$ such that
\begin{equation}
  w_-\leq w(s)\leq w_+ \qquad\text{for almost every }s.
  \label{eq:weight}
\end{equation}
Let $\Omega,\sigma\in L^2(0,S;\RR)$ be fixed.
We use a fixed framing and endpoint convention in which an admissible real
phase lift $\chi\in H^1(0,S;\RR)$ satisfies
\begin{equation}
  \chi(S)-\chi(0)=2\pi n.
  \label{eq:winding}
\end{equation}
The phase $e^{i\chi}$ is therefore periodic.  Define
\begin{equation}
  q(s):=\chi'(s)-m\Omega(s),
  \qquad
  \E[\chi]:=\int_0^S w(s)\bigl(q(s)-\sigma(s)\bigr)^2\dd s.
  \label{eq:functional}
\end{equation}
Introduce the fixed quantities
\begin{equation}
  \Phi:=\int_0^S\Omega(s)\dd s,
  \qquad
  J:=\int_0^S\sigma(s)\dd s,
  \qquad
  I:=\int_0^S\frac{\dd s}{w(s)}.
  \label{eq:definitions}
\end{equation}
By \eqref{eq:weight}, $0<I<\infty$.  The admissible space in winding
sector $n$ is
\begin{equation}
  \mathcal{A}_n:=\bigl\{\chi\in H^1(0,S;\RR):
                 \chi(S)-\chi(0)=2\pi n\bigr\}.
\end{equation}
All geometric data, the framing, $m$, and $n$ are held fixed during
minimization.

\section{Exact minimum and sharp remainder}

\begin{theorem}[Exact weighted phase-locking minimum]
\label{thm:main}
Under the assumptions above, set
\begin{equation}
  \Delta_n:=2\pi n-m\Phi-J.
  \label{eq:delta}
\end{equation}
Then
\begin{equation}
  \boxed{\displaystyle
  \min_{\chi\in\mathcal{A}_n}\E[\chi]
  =\frac{\Delta_n^2}{I}
  =\frac{(2\pi n-m\Phi-J)^2}{\displaystyle\int_0^S w^{-1}\dd s}.}
  \label{eq:minimum}
\end{equation}
The minimizing covariant phase gradient is, uniquely up to equality almost
everywhere,
\begin{equation}
  \boxed{\displaystyle
  q_*(s)=\sigma(s)+\frac{\Delta_n}{I\,w(s)}.}
  \label{eq:qstar}
\end{equation}
An admissible minimizing phase is
\begin{equation}
  \boxed{\displaystyle
  \chi_*(s)=\chi_0+\int_0^s
  \left[m\Omega(t)+\sigma(t)+\frac{\Delta_n}{I\,w(t)}\right]\dd t,}
  \label{eq:chistar}
\end{equation}
where $\chi_0\in\RR$ is arbitrary.  Every minimizing phase differs from
$\chi_*$ by an additive constant.  Furthermore, every
$\chi\in\mathcal{A}_n$ satisfies the exact identity
\begin{equation}
  \boxed{\displaystyle
  \E[\chi]
  =\frac{\Delta_n^2}{I}
   +\int_0^S w(s)\lvert q(s)-q_*(s)\rvert^2\dd s.}
  \label{eq:remainder}
\end{equation}
\end{theorem}

\begin{proof}
Integrating the definition of $q$ and using the winding constraint gives
\begin{equation}
  \int_0^S q\dd s
  =\chi(S)-\chi(0)-m\int_0^S\Omega\dd s
  =2\pi n-m\Phi.
  \label{eq:qconstraint}
\end{equation}
Hence for the mismatch $f:=q-\sigma$,
\begin{equation}
  \int_0^S f\dd s=2\pi n-m\Phi-J=\Delta_n.
  \label{eq:fconstraint}
\end{equation}
By the Cauchy--Schwarz inequality,
\begin{align}
  \Delta_n^2
  &=\left(\int_0^S
     \sqrt{w(s)}f(s)\frac{1}{\sqrt{w(s)}}\dd s\right)^2\\
  &\leq\left(\int_0^S w(s)f(s)^2\dd s\right)
          \left(\int_0^S\frac{\dd s}{w(s)}\right)
   =I\E[\chi].
\end{align}
Because $I>0$, this yields
\begin{equation}
  \E[\chi]\geq\frac{\Delta_n^2}{I}.
  \label{eq:lower}
\end{equation}
Cauchy--Schwarz is an equality precisely when
$\sqrt{w}\,f=C/\sqrt{w}$ almost everywhere for some constant $C$.
Thus $f=C/w$, and \eqref{eq:fconstraint} implies
\begin{equation}
  \Delta_n=C\int_0^S w^{-1}\dd s=CI,
  \qquad C=\frac{\Delta_n}{I}.
\end{equation}
The only possible minimizing covariant gradient is consequently
\eqref{eq:qstar}.

To prove that the lower bound is attained, define $\chi_*$ by
\eqref{eq:chistar}.  The stated regularity assumptions ensure
$\chi_*\in H^1(0,S)$, and
\begin{align}
  \chi_*(S)-\chi_*(0)
    &=m\Phi+J+\frac{\Delta_n}{I}
      \int_0^S\frac{\dd s}{w(s)}\\
    &=m\Phi+J+\Delta_n=2\pi n.
\end{align}
Thus $\chi_*\in\mathcal{A}_n$, and $q_*=\chi_*'-m\Omega$.
Its energy is
\begin{equation}
  \E[\chi_*]
  =\int_0^S w\left(\frac{\Delta_n}{Iw}\right)^2\dd s
  =\frac{\Delta_n^2}{I^2}\int_0^S w^{-1}\dd s
  =\frac{\Delta_n^2}{I}.
\end{equation}
This proves the exact minimum.

For the remainder identity, let $h:=q-q_*$.  Since $\chi$ and $\chi_*$
have the same winding number,
\begin{equation}
  \int_0^S h\dd s
  =\int_0^S(\chi'-\chi_*')\dd s=0.
  \label{eq:hzero}
\end{equation}
Also, $w(q_*-\sigma)=\Delta_n/I$.  Expanding the square yields
\begin{align}
 \E[\chi]
 &=\int_0^S w\bigl[(q_*-\sigma)+h\bigr]^2\dd s\\
 &=\E[\chi_*]
   +\frac{2\Delta_n}{I}\int_0^S h\dd s
   +\int_0^S wh^2\dd s\\
 &=\frac{\Delta_n^2}{I}+\int_0^S w\lvert q-q_*\rvert^2\dd s.
\end{align}
Positivity of $w$ shows that equality holds only for $q=q_*$ almost
everywhere.  In that case $(\chi-\chi_*)'=0$ almost everywhere,
so $\chi-\chi_*$ is constant.  This proves the asserted uniqueness.
\end{proof}

\section{Consequences}

\begin{corollary}[Criterion for an exact lock]
\label{cor:lock}
There exists an admissible phase with $q=\sigma$ almost everywhere if and
only if
\begin{equation}
  \boxed{2\pi n=m\Phi+J.}
\end{equation}
Equivalently, $\min_{\mathcal{A}_n}\E=0$ if and only if $\Delta_n=0$.
\end{corollary}
\begin{proof}
The conclusion follows from \eqref{eq:minimum}, or directly by integrating
$q=\sigma$ in \eqref{eq:qconstraint}.
\end{proof}

\begin{corollary}[Strict phase-sector variational stability]
\label{cor:stability}
Let $\eta\in H^1(0,S;\RR)$ have periodic endpoint values
$\eta(S)=\eta(0)$.  Then for every $\varepsilon\in\RR$,
\begin{equation}
  \E[\chi_*+\varepsilon\eta]
  =\E[\chi_*]
   +\varepsilon^2\int_0^S w(s)\lvert\eta'(s)\rvert^2\dd s.
  \label{eq:perturb}
\end{equation}
In particular,
\begin{equation}
  \left.\frac{\dd^2}{\dd\varepsilon^2}
      \E[\chi_*+\varepsilon\eta]\right|_{\varepsilon=0}
  =2\int_0^S w\lvert\eta'\rvert^2\dd s\geq0,
\end{equation}
with equality if and only if $\eta$ is constant almost everywhere.
\end{corollary}
\begin{proof}
The perturbation preserves the winding constraint, and its covariant
gradient differs from $q_*$ by $\varepsilon\eta'$.  Apply
\eqref{eq:remainder}.
\end{proof}

\begin{remark}[Meaning of stability]
Corollary~\ref{cor:stability} is a statement about a fixed-geometry,
phase-only variational functional.  It does not assert dynamical, orbital,
or spectral stability for a full complex-field evolution.  In particular,
it does not allow the winding number to change during the perturbation.
\end{remark}

\begin{proposition}[Conditional comparison of winding sectors]
\label{prop:sectors}
Suppose the geometric data $\GG,w,\sigma,\Omega$ and integer $m$ are fixed,
all integer winding sectors are admissible, and the complete
sector-dependent energy is $\E$ plus a term independent of $n$.
If the system is allowed to compare the global minimum energies of these
sectors, define
\begin{equation}
  \zeta:=\frac{m\Phi+J}{2\pi}.
\end{equation}
Then the integer sectors of least energy are exactly
\begin{equation}
  \boxed{\displaystyle
  \mathcal{N}_{\min}
    =\operatorname*{arg\,min}_{n\in\ZZ}|n-\zeta|.}
\end{equation}
There is one minimizing integer unless $\zeta\in\ZZ+\tfrac12$,
when two adjacent integers are tied.
\end{proposition}
\begin{proof}
Theorem~\ref{thm:main} gives
\begin{equation}
  \E_{\min}(n)
  =\frac{(2\pi n-m\Phi-J)^2}{I}
  =\frac{4\pi^2}{I}(n-\zeta)^2.
\end{equation}
The prefactor is positive and independent of $n$, so minimizing over
integers is equivalent to selecting the integer(s) nearest to $\zeta$.
\end{proof}

\begin{remark}[Limits of inter-sector selection]
Proposition~\ref{prop:sectors} compares minimum costs; it does not establish
that the physical dynamics can change $n$.  With a continuous nonvanishing
complex field on a fixed closed contour, winding is topologically conserved
under continuous evolution.  A change of winding requires additional
conditions or dynamics, for example a field zero (phase slip), altered
boundary data, or a change of the contour topology.  Moreover, omitted
$n$-dependent energies would change the sector comparison.
\end{remark}

\begin{remark}[Frame conventions]
The constraint $\int q\dd s=2\pi n-m\Phi$ assumes the stated phase-lift
and framing convention.  If the normal frame has a nontrivial seam
identification, its transition phase must be incorporated consistently into
the endpoint condition and holonomy.  The proof then applies with the
correspondingly adjusted, gauge-consistent total closure mismatch.
\end{remark}

\section{Uniform-weight example}

If $w(s)\equiv w_0>0$, $\Omega\equiv0$, and $\sigma(s)\equiv\sigma_0$,
then $I=S/w_0$, $J=S\sigma_0$, and
\begin{equation}
  \E_{\min}(n)=\frac{w_0}{S}(2\pi n-S\sigma_0)^2,
  \qquad
  q_*(s)=\frac{2\pi n}{S}.
\end{equation}
Thus the optimal phase gradient is uniform, but it need not coincide with
the preferred local rate $\sigma_0$.  Exact locking occurs precisely when
$S\sigma_0=2\pi n$.

\section*{Scope}
The theorem is an exact result for the specified weighted phase-only
functional, closed-loop geometry, winding sector, and framing.  Its proof
uses the sharp weighted Cauchy--Schwarz inequality; it does not provide an
independent physical derivation of the locking energy, a mechanism for
transitions between winding sectors, or a prediction for geometry or
spectral frequencies.

\noindent\textbf{Supporting analysis:}
\url{https://github.com/rickyjreyes/geometry_of_resonance/blob/main/docs/WINDING_SELECTION_ANALYSIS.md}

\end{document}
